\subsection{Performance and Optimization}\label{subsec:num_performance}

This section addresses computational performance considerations for large-scale numerical integrations.

\subsubsection{Vectorization and SIMD}

Modern processors execute multiple operations in parallel (SIMD). Exploit this:

\begin{itemize}
    \item **Outer loop over $\ell$ values**: Compute all $C_\ell(\ell=0,1,2,...)$ simultaneously
    \item **Batch kernel evaluations**: Evaluate $j_\ell(kr)$ for many $(k, r)$ pairs at once
    \item **Avoid scalar loops**: Use array operations (map, reduce) instead of for-loops
\end{itemize}

**Julia tools**: `@turbo`, `@tturbo` (LoopVectorization.jl), `@tullio` (Tullio.jl)

\subsubsection{Memory access patterns}

\begin{itemize}
    \item **Sequential access**: Much faster than random access (CPU cache efficiency)
    \item **Cache-friendly layout**: Store data as SoA (structure of arrays) not AoS (array of structures)
    \item **Pre-allocation**: Avoid repeated memory allocation inside tight loops
\end{itemize}

\subsubsection{Bessel function caching}

Spherical Bessel functions are expensive. Cache when possible:

\begin{itemize}
    \item If evaluating $j_\ell(kr_i)$ for multiple $i$ at fixed $\ell, k$: compute once, reuse
    \item Use recurrence relations: $j_{\ell+1}(x) = \frac{(2\ell+1) j_\ell(x) - x j_{\ell-1}(x)}{x}$
    \item Pre-tabulate for common $\ell$ values
\end{itemize}

\subsubsection{Convergence vs accuracy trade-off}

\begin{itemize}
    \item **Tight tolerance** ($\text{rtol} = 10^{-12}$): More accurate but 10× slower
    \item **Loose tolerance** ($\text{rtol} = 10^{-6}$): Fast but risky for power spectrum (convolutions amplify errors)
    \item **Compromise**: $\text{rtol} = 10^{-9}, \text{atol} = 10^{-12}$ (often sufficient)
\end{itemize}

\remark{
For cosmological predictions competing with observation (1--5\% precision), target ~0.1\% accuracy in integrals to avoid bias.
}

See \cref{subsec:num_benchmarking} for performance measurement techniques.
